\documentclass[ecta,nameyear,final,supplement]{econsocart}

\usepackage{amsmath,amssymb,bm}
\usepackage{booktabs}
\usepackage{array}
\usepackage{tabularx}

\startlocaldefs
\numberwithin{equation}{section}
\renewcommand{\theequation}{S\arabic{section}.\arabic{equation}}
\renewcommand{\thesection}{S\arabic{section}}
\newcommand{\Prob}{\mathbb{P}}
\newcommand{\E}{\mathbb{E}}
\newcommand{\Ind}{\mathbf{1}}
\newcommand{\calA}{\mathcal{A}}
\newcommand{\ThetaK}{\Theta_K}
\newcommand{\thetak}{\theta_K}
\newcommand{\lambdat}{\widetilde{\lambda}}
\newcommand{\LH}{\mathcal{L}^{H}}
\newcommand{\LD}{\mathcal{L}^{d}}
\newcommand{\LC}{\mathcal{L}^{C}}
\newcommand{\LI}{\mathcal{L}^{I}}
\newcommand{\LT}{\mathcal{L}}
\newcommand{\argmin}{\operatorname*{arg\,min}}
\newcommand{\argmax}{\operatorname*{arg\,max}}
\endlocaldefs

\begin{document}

\begin{frontmatter}
\title{Supplement to ``Identifying Rational Types in Unknown Environments''\\
Appendix 4: Streamlit Calibration Model}
\runtitle{Streamlit Calibration Supplement}

\begin{aug}
\author[add1,add2]{\fnms{Humberto}~\snm{Bernal}}
\address[add1]{%
\orgdiv{Economics Program},
\orgname{Universidad Colegio Mayor de Cundinamarca}}
\address[add2]{%
\orgdiv{Ph.D. in Economics},
\orgname{Universidad de los Andes}}
\ead[label=e1]{hbernal@uniandes.edu.co}
\ead[label=e2]{hbernalc@universidadmayor.edu.co}
\end{aug}
\end{frontmatter}

\section{Scope of This Computational Appendix}

This appendix documents only the equations that are used by the Streamlit
calibration app in \path{../Streamlit/app.py}. It follows the order of the main
manuscript: first the empirical objects that discipline priors and hazards,
then the theoretical mechanism implemented in the app, then the feasibility
diagnostics, and finally the way in which these equations generate Table 5.2.

The macroeconomic VEC and LLT estimates in the main paper are not reproduced
here because the Streamlit app does not re-estimate them and does not use them
directly to compute Table 5.2. The equations below are the active estimated or
calibrated objects used by the app.

\section{Empirical Inputs Used by the App}

\subsection{Outcome Model and Prior Over Types}

The first empirical input is the multinomial outcome model from the manuscript.
With payment as the reference outcome, the estimated equation is
\begin{equation}
\ln\!\left(
\frac{\Pr(Y_{i,m}=j\mid \mathbf X_{i,m})}
{\Pr(Y_{i,m}=1\mid \mathbf X_{i,m})}
\right)
=
\alpha_j+\mathbf G_i'\bm\beta_{1,j}
+\mathbf V_i'\bm\beta_{2,j}
+\mathbf C_{i,m}'\bm\beta_{3,j}
+\mathbf T_i'\bm\beta_{4,j},
\quad j\in\{0,2,3\}.
\label{eq:app-mnl}
\end{equation}
In the app, this block is used to discipline the initial type probabilities and
case covariate adjustments. The prior used by the Streamlit state is a softmax:
\begin{equation}
\mu_0(\theta\mid z,\theta_V)
=
\frac{\exp(\varpi_\theta+\eta_{\theta,z}+\xi_{\theta,\theta_V})}
{\sum_{\theta'\in\ThetaK}
\exp(\varpi_{\theta'}+\eta_{\theta',z}+\xi_{\theta',\theta_V})}.
\label{eq:app-prior}
\end{equation}
The inputs are the type baseline $\varpi_\theta$, the region adjustment
$\eta_{\theta,z}$, and the victim-profile adjustment $\xi_{\theta,\theta_V}$.
The output is the initial posterior vector $\mu_0$ used by Table 5.2.

\subsection{Cause-Specific Duration Model}

The second empirical input is the Cox competing-risk block. For each cause
$j$, the manuscript estimates
\begin{equation}
h_j(t\mid\mathbf X)
=
h_{0j}(t)\exp(\mathbf X'\widetilde{\bm\beta}_j),
\qquad
S_j(t\mid\mathbf X)
=
\exp\!\left[-H_{0j}(t)\exp(\mathbf X'\widetilde{\bm\beta}_j)\right].
\label{eq:app-cox}
\end{equation}
The app uses this information through baseline risks $\lambda_{j0}$ and
type-cause coefficients $\beta_{K,j}(\theta)$, stored in the model object and
updated by the calibration interface. These coefficients feed the structural
intensities in Section~\ref{sec:app-risk}.

\section{Theoretical Model Implemented in Streamlit}

\subsection{State, History, and Cycle}

The app implements the paper's dynamic history as the period-by-period record
of instruments, executed actions, physical outcomes, and detection signals:
\begin{gather}
h_t=
\left(
t,\alpha^\ast_{[0,t-1]},\gamma^\ast_{[0,t-1]},
\widetilde a^F_{[0,t-1]},\widetilde a^K_{[0,t-1]},
\widetilde a^S_{[0,t-1]},m_{[0,t-1]},d_{[0,t-1]}
\right),\label{eq:app-history}\\
h_{t+1}=
\left(h_t,\alpha_t^\ast,\gamma_t^\ast,
\widetilde a_t^F,\widetilde a_t^K,\widetilde a_t^S,m_t,d_t\right).
\label{eq:app-history-update}
\end{gather}
The terminal event is
\begin{equation}
\tau=\inf\{t\ge0:m_t\in\mathcal M_{\mathrm{term}}\}.
\label{eq:app-stopping}
\end{equation}
When Table 5.2 runs in realized-draw mode, the simulated event stops after the
first $m_t\ne\mathrm{cont}$; the diagnostic tables still retain the probability
vector used to generate the draw.

\subsection{MDG Implementation Kernel}

For each agent $i\in\{K,S,F\}$, the app maps the intended action
$a_t^{i\ast}$ into an executed action $\widetilde a_t^i$ with the MDG logit
kernel:
\begin{equation}
\Prob_{\mathrm I,i}
\left(\widetilde a_t^i=a\mid a_t^{i\ast},X_t\right)
=
\frac{\exp\{\Ind(a=a_t^{i\ast})/T_t\}}
{\sum_{a'\in\calA^i}\exp\{\Ind(a'=a_t^{i\ast})/T_t\}}.
\label{eq:app-mdg}
\end{equation}
The temperature used in the app is
\begin{equation}
T_t
=
T_0\frac{H(\mu_t)}{H(\mu_0)}
\max\{\exp(-\eta_{\mathrm{cal}}t),\bar c\},
\qquad
H(\mu_t)=-\sum_{\theta\in\ThetaK}\mu_t(\theta)\ln\mu_t(\theta).
\label{eq:app-temperature}
\end{equation}
The app also reports the bounded operational noise
\begin{equation}
\varepsilon_t=\min\{0.5,\max\{0,\varepsilon_0T_t\}\}.
\label{eq:app-mdg-eps}
\end{equation}
For Table 5.2, the implemented-action likelihood is
\begin{equation}
\LI_t(\theta)
=
\prod_{i\in\{K,S,F\}}
\Prob_{\mathrm I,i}
\left(\widetilde a_t^i\mid a_t^{i\ast},X_t\right),
\label{eq:app-li}
\end{equation}
computed by the implementation-likelihood block in \path{app.py}.

\subsection{Competing-Risk Technology}
\label{sec:app-risk}

The app evaluates the structural proportional-risk equation
\begin{equation}
\lambda_j(t\mid\mathcal C_t,\theta)
=
\lambda_{j0}(t)\exp(x_t'\beta_j(\theta)),
\qquad j\in\{1,2,3\}.
\label{eq:app-lambda-main}
\end{equation}
Operationally, the app uses
$\lambdat_j(t\mid\thetak)=M(t)\lambda_{j0}\exp(E_j)$, with
\begingroup\small
\begin{align}
E_1
&=
\beta_{K,1}(\thetak)+\beta_{z,1}+\beta_{F,1}-\beta_{V,1}
+\beta_{S,1}
-\zeta_{\alpha,1}\alpha_t-\zeta_{\gamma,1}\gamma_t
-\zeta_{d,1}p_{\mathrm{det},t}
+\phi_{F,1}\Ind_F+\phi_{K,1}^{cont}\Ind_K^{cont},
\label{eq:app-exp-pay}\\
E_2
&=
\beta_{K,2}(\thetak)+\beta_{z,2}+\beta_{S,2}
+\zeta_{\alpha,2}\alpha_t+\zeta_{\gamma,2}\gamma_t
-\zeta_{d,2}p_{\mathrm{det},t}
-\phi_{F,2}\Ind_F+\phi_{K,2}^{kill}\Ind_K^{kill}
+\phi_{K,2}^{cont}\Ind_K^{cont},
\label{eq:app-exp-kill}\\
E_3
&=
\beta_{K,3}(\thetak)+\beta_{z,3}-\beta_{S,3}
+\zeta_{\alpha,3}\alpha_t+\zeta_{\gamma,3}\gamma_t
+\zeta_{d,3}p_{\mathrm{det},t}
+\zeta_R\Ind_S^{res}
-\phi_{F,3}\Ind_F+\phi_{K,3}^{cont}\Ind_K^{cont}.
\label{eq:app-exp-res}
\end{align}
\endgroup
Equivalently,
\begin{equation}
\lambdat_j(t\mid\thetak)=M(t)\lambda_{j0}\exp(E_j),
\qquad j=1,2,3.
\label{eq:app-lambda-strategic}
\end{equation}
Here
\begin{equation}
M(t)=\min\{1,(t/T_{\mathrm{mad}})^2\},
\qquad
\lambdat_4(t)=\lambda_4.
\label{eq:app-maturity-lambda4}
\end{equation}
The daily probabilities used by Table 5.2 are computed as
\begin{align}
p_{\mathrm{Cont},t\mid\thetak}
&=
\exp\!\left[-\Delta t\sum_{j=1}^{4}\lambdat_j(t\mid\thetak)\right],
\label{eq:app-pcont}\\
q_t(\thetak)&=1-p_{\mathrm{Cont},t\mid\thetak},
\label{eq:app-q}\\
\xi_j(t\mid\thetak)
&=
\frac{\lambdat_j(t\mid\thetak)}
{\sum_{\ell=1}^{4}\lambdat_\ell(t\mid\thetak)},
\qquad
h_j(t\mid\thetak)=q_t(\thetak)\xi_j(t\mid\thetak).
\label{eq:app-xi-h}
\end{align}
These calculations are performed by \path{mechanism_competitive_hazards_at_t}.
The returned vector is
\[
\left(
p_{\mathrm{Cont}},h_{\mathrm{pay}},h_{\mathrm{kill}},
h_{\mathrm{res}},h_{\mathrm{rel}}
\right).
\]

\subsection{Detection and Communication Signals}

The detection probability used by the app is type-specific in its intercept:
\begin{equation}
p_{\mathrm{det},t}(\thetak)
=
\Lambda\!\left(\eta_0(\thetak)+\eta_1\alpha_t^\ast+\eta_2\gamma_t^\ast\right),
\qquad
\Lambda(u)=\frac{1}{1+\exp(-u)}.
\label{eq:app-detection}
\end{equation}
The detection likelihood is Bernoulli:
\begin{equation}
\LD_t(d_t\mid\thetak)
=
p_{\mathrm{det},t}(\thetak)^{d_t}
\left[1-p_{\mathrm{det},t}(\thetak)\right]^{1-d_t}.
\label{eq:app-ld}
\end{equation}
When the communication block is active, the app evaluates either a voice
likelihood or a silence likelihood. The voice signal is
\begin{equation}
x_t^{obs}
=
x_t^{true}(\thetak)+\varepsilon_L+\varepsilon_S,
\qquad
\varepsilon_L\sim N(0,\Sigma_L),\quad
\varepsilon_S\sim N(0,\Sigma_S),
\label{eq:app-voice}
\end{equation}
with diagonal likelihood
\begin{equation}
\mathcal L_{\mathrm{voz},t}(\thetak)
\propto
\exp\!\left[-\frac12
\sum_k
\left(\frac{x_{t,k}^{obs}-m_k(\thetak)}{\sigma_k(\thetak)}\right)^2
\right].
\label{eq:app-voice-likelihood}
\end{equation}
If no call is observed, the app uses
\begin{equation}
\LC_t(\thetak\mid V_t=0)
=
\left[1-\pi_{\mathrm{call}}(\thetak)\right]^{\omega_{\mathrm{voz}}}.
\label{eq:app-silence}
\end{equation}

\subsection{Physical Likelihood and Posterior}

Given the competing-risk vector, the physical likelihood is
\begin{equation}
\LH_t(m_t\mid\thetak)
=
p_{\mathrm{Cont},t\mid\thetak}^{\Ind\{m_t=\mathrm{cont}\}}
\prod_{j=1}^{4}
h_j(t\mid\thetak)^{\Ind\{m_t=g(j)\}},
\label{eq:app-lh}
\end{equation}
where $g(1)=\mathrm{pay}$, $g(2)=\mathrm{kill}$,
$g(3)=\mathrm{res}$, and $g(4)=\mathrm{rel}$. The total likelihood used for
Table 5.2 is
\begin{equation}
\LT_t(\thetak)
=
\LI_t(\thetak)\LH_t(m_t\mid\thetak)
\LD_t(d_t\mid\thetak)\LC_t(\thetak).
\label{eq:app-total-likelihood}
\end{equation}
The posterior update is
\begin{equation}
\mu_{t+1}(\thetak)
=
\frac{\mu_t(\thetak)\LT_t(\thetak)}
{\sum_{\theta'\in\ThetaK}\mu_t(\theta')\LT_t(\theta')}.
\label{eq:app-bayes}
\end{equation}
This is computed by \path{build_t0_bayesian_posterior_report} and
\path{bayesian_posterior_update}. The output is the likelihood decomposition
and the new belief vector used in the next cycle.

\section{Policy and Feasibility Equations Used by the App}

\subsection{State Objective}

The app computes the state's branch loss as
\begin{equation}
L_t(a_t^S,\alpha_t,\gamma_t,\thetak,\iota_t)
=
\begin{cases}
V_t^R(\iota_t,\widehat\theta_t,\thetak,\alpha_t,\gamma_t),
&a_t^S=a_{\mathrm{res}},\\
V_t^N(\thetak,\alpha_t,\gamma_t),
&a_t^S=a_{\mathrm{neg}}.
\end{cases}
\label{eq:app-state-loss}
\end{equation}
The numerical optimizer for Table 5.2 solves
\begin{equation}
\begin{aligned}
(a_t^{S\ast},\alpha_t^\ast,\gamma_t^\ast)
&\in
\argmin_{\Gamma_t(\mu_t)}
\Biggl[
\sum_{\theta\in\ThetaK}\mu_t(\theta)
L_t(a_t^S,\alpha_t,\gamma_t,\theta,\iota_t)\\
&\qquad
+\Pi_{t,a^S}^S(\alpha_t,\gamma_t;\mu_t)
-\psi_H\Delta H_t(\alpha_t,\gamma_t)
\Biggr].
\end{aligned}
\label{eq:app-state-objective}
\end{equation}
The penalty around the Bayesian benchmark and the entropy gain are
\begin{align}
\Pi_{t,b}^S(\alpha_t,\gamma_t;\mu_t)
&=
\chi_\alpha(\alpha_t-\alpha_{t,b}^{\mu})^2
+\chi_\gamma(\gamma_t-\gamma_{t,b}^{\mu})^2,
\label{eq:app-deviation}\\
\Delta H_t(\alpha_t,\gamma_t)
&=
H(\mu_t)
-\sum_m
\Prob(m_t=m\mid\mu_t,\alpha_t,\gamma_t)H(\mu_{t+1}^m).
\label{eq:app-deltaH}
\end{align}
The app compares the rescue and negotiation floors and selects rescue when
$\widetilde V_t^{R,\ast}\le \widetilde V_t^{N,\ast}$.

\subsection{Private-Agent Equations and Diagnostics}

The app evaluates the captor's continuation cost as
\begin{equation}
C_t(\gamma_t,\thetak)
=
\phi(\thetak)\exp\{\kappa_c(\thetak)\gamma_t\}+\nu(\thetak).
\label{eq:app-k-cost}
\end{equation}
The family compares cooperation with clandestine payment:
\begin{align}
\mathcal U_t^F(a_{\mathrm{coop}})
&=
\E_{\theta\mid\mathcal I_t^F}
\left[\widetilde p_{\mathrm{surv},t}(\theta)\right]V_L
-e_t(\gamma_t,\theta_F),
\label{eq:app-f-coop}\\
\mathcal U_t^F(a_{\mathrm{col}})
&=
\E_{\theta\mid\mathcal I_t^F}
\left[\widetilde p_{\mathrm{rel},t}(\theta)\right]V_L
-R-p_{\mathrm{det},t}F_{\mathrm{col}}.
\label{eq:app-f-col}
\end{align}
The feasibility diagnostics shown with Table 5.2 and Tab 6 are
\begin{align}
IC_t^K&:
\quad
\E_{\theta\sim\mu_t}
\left[
V_t^K(a^\ast(\theta)\mid\theta,\alpha_t,\gamma_t)
-V_t^K(a^\ast(\theta_j)\mid\theta,\alpha_t,\gamma_t)
\right]\ge0,
\label{eq:app-ic-k}\\
IR_t^K&:
\quad
\E_{\theta\sim\mu_t}
\left[
U^K_{\mathrm{rel}}(\theta)
-\max\{V^K_{\mathrm{cont},t}(\theta,\alpha_t,\gamma_t),
U^K_{\mathrm{kill}}(\theta,\theta_S)\}
\right]\ge0,
\label{eq:app-ir-k}\\
IR_t^F&:
\quad
\mathcal U_t^F(a_{\mathrm{coop}})
\ge
\mathcal U_t^F(a_{\mathrm{col}}).
\label{eq:app-ir-f}
\end{align}

\section{How the App Obtains Table 5.2}

Table 5.2 is produced by iterating the equations above. The sequence implemented
in \path{app.py} is:
\begin{equation}
\mu_t
\rightarrow
(a_t^{K\ast},a_t^{S\ast},a_t^{F\ast})
\rightarrow
\widetilde A_t
\rightarrow
(\lambdat_j,p_{\mathrm{Cont}},h_j,p_{\mathrm{det}})
\rightarrow
(m_t,d_t,V_t,x_t^{obs})
\rightarrow
\mu_{t+1}.
\label{eq:app-cycle}
\end{equation}

At $\tau=0$, the app starts from $\mu_0$ in~\eqref{eq:app-prior}. It computes
the state policy $(a_0^{S\ast},\alpha_0^\ast,\gamma_0^\ast)$ from
\eqref{eq:app-state-objective}, computes the private intended actions, applies
the MDG implementation kernel~\eqref{eq:app-mdg}, evaluates
$p_{\mathrm{det}}$ using~\eqref{eq:app-detection}, computes the physical
probabilities through~\eqref{eq:app-pcont}--\eqref{eq:app-xi-h}, and then forms
the likelihood~\eqref{eq:app-total-likelihood}. The posterior column
$\mu_1$ is obtained from~\eqref{eq:app-bayes}.

For each subsequent cycle $\tau\ge1$, the posterior inherited from the previous
cycle becomes the new prior. The same optimizer, implementation kernel,
hazard block, detection block, likelihood block, and Bayes rule are applied
again. Thus, every column of Table 5.2 is a repeated application of
\eqref{eq:app-cycle}.

\begin{table}[t]
\centering
\caption{Equation-to-output map for Table 5.2}
\begin{tabularx}{\textwidth}{>{\raggedright\arraybackslash}p{0.32\textwidth}X}
\toprule
Table 5.2 object & Equation used by the app \\
\midrule
Prior and posterior beliefs & Prior softmax~\eqref{eq:app-prior} and Bayes
rule~\eqref{eq:app-bayes}. \\
Optimal state instruments & State objective~\eqref{eq:app-state-objective},
penalty~\eqref{eq:app-deviation}, and entropy gain~\eqref{eq:app-deltaH}. \\
MDG intended/executed actions & Kernel~\eqref{eq:app-mdg} and implementation
likelihood~\eqref{eq:app-li}. \\
Physical outcome probabilities & Intensities~\eqref{eq:app-lambda-strategic}
and probabilities~\eqref{eq:app-pcont}--
\eqref{eq:app-xi-h}. \\
Detection signal & Logistic detection~\eqref{eq:app-detection} and Bernoulli
likelihood~\eqref{eq:app-ld}. \\
Posterior decomposition & Physical likelihood~\eqref{eq:app-lh}, total
likelihood~\eqref{eq:app-total-likelihood}, and normalization in
\eqref{eq:app-bayes}. \\
Feasibility diagnostics & $IC^K$, $IR^K$, and $IR^F$ in
\eqref{eq:app-ic-k}--\eqref{eq:app-ir-f}. \\
\bottomrule
\end{tabularx}
\end{table}

The resulting objects are stored by the app in the dynamic Table 5.2 session
state, including \path{dynamic_cycles52}, \path{dynamic_cycles_diag52},
\path{dynamic_cycles_stop52}, \path{first_cycle_table52}, and
\path{first_cycle_post54}. Tab 6 reads these same objects; it does not change
the calculations used to produce Table 5.2.

\end{document}
